% Section 1.3: Contraction: the canonical pairing | Part: 1 | Depends on: none | Status: done | Original lines: 453-501
%!TEX root = ../../../main.tex
\subsection{Contraction: the canonical pairing}\label{sec:contraction}

The evaluation of a covector on a vector,
\begin{equation}\label{eq:pairing}
\langle\omega,v\rangle=\omega(v)=\omega_i\,v^i\qquad(\text{sum over }i),
\end{equation}
is called the \emph{contraction} of $\omega\otimes v$: the index $i$ is ``contracted'' between the lower slot of $\omega$ and the
upper slot of $v$. It needs \emph{no metric, no basis and no inner product}: it is part of the very definition of $V^*$.
Let us see, step by step, why its formula in components is $\omega_iv^i$:
\begin{align*}
\omega(v)&=\omega\bigl(v^je_j\bigr) &&\text{(expand $v$ in the basis $e_j$)}\\
&=v^j\,\omega(e_j) &&\text{(linearity of $\omega$ in its argument)}\\
&=v^j\,\omega_i\,e^i(e_j) &&\text{(expand $\omega=\omega_ie^i$)}\\
&=v^j\,\omega_i\,\delta^i{}_j &&\text{(duality $e^i(e_j)=\delta^i{}_j$)}\\
&=\omega_i\,v^i &&\text{($\delta^i{}_j$ keeps only the term $j=i$).}
\end{align*}
\numcheck In $\R^2$ let $\omega=3\,\dd x-2\,\dd y$ (so $\omega_x=3$, $\omega_y=-2$) and $v=2\partial_x+5\partial_y$
(so $v^x=2$, $v^y=5$). Then $\omega(v)=3\cdot2+(-2)\cdot5=-4$.

% Figure | used in parts/01-algebra/01-vectors-covectors/03-contraction.tex
%!TEX root = ../../main.tex
\begin{figure}[ht]
\centering
\begin{tikzpicture}[>=Stealth]
  \foreach \i in {0,...,7}{\draw[blue!65!black,thin] (\i*0.62,-0.45)--(\i*0.62+0.42,1.65);}
  \draw[->,red!75!black,very thick] (0.12,0.15)--(4.75,1.20);
  \node[red!75!black,below right] at (2.6,0.62) {$v$};
  \node[blue!65!black,above] at (0.15,1.65) {$\omega$};
  \node[align=left,anchor=west] at (5.9,0.65) {number of sheets pierced\\[2pt] $\Rightarrow\ \omega(v)\approx 3.9$};
\end{tikzpicture}
\caption{Geometric picture of contraction. A covector $\omega$ is a ``stack of parallel sheets'' (its level sets); the vector
$v$ is an arrow; $\omega(v)$ counts how many sheets the arrow pierces. Doubling $v$ doubles the count, and reversing $v$
changes the sign: this is linearity.}
\label{fig:sheets}
\end{figure}


\begin{keyidea}
\textbf{Contraction is the metric-free pairing $V^*\times V\to\R$.} Everything later that looks like a ``dot product'' in physics
-- work $\mathbf F\cdot\dd\mathbf r$, power, flux, the pairing of a momentum with a symmetry -- is, when examined, a contraction
$\langle\omega,v\rangle$ of a covector with a vector. A metric is needed only when one wants to turn an arrow into a ruler or
vice versa (Section~\ref{sec:metric}).
\end{keyidea}

Contraction is the first member of a family: it extends to tensors of any type (Section~\ref{sec:tensors}), to
$k$-forms as the \emph{interior product} $\iota_v\alpha$ (Section~\ref{sec:interior}), and to multivectors as the
\emph{left contraction} $a\lc A$ of geometric algebra (Section~\ref{sec:ops}); Section~\ref{sec:synth1} compares these four
guises. On a manifold the same pairing is performed at every point between the cotangent space and the tangent space
(Section~\ref{sec:tangent}), and integrating along a curve will turn out to be nothing more than contracting a $1$-form with the
velocity (Section~\ref{sec:int1}).
